\documentclass[acmtog,anonymous,review]{acmart}

\acmSubmissionID{TBD}

\settopmatter{printfolios=true}
\renewcommand\shortauthors{Anonymous Submission}

\begin{document}

\title{Geometry-Locked Multispectral Gaussian Splatting for UAV Imaging: Supplementary Material}

\author{Anonymous Author(s)}
\affiliation{%
  \institution{Anonymous Institution}
  \country{Anonymous}
}

\ccsdesc[500]{Computing methodologies~Image-based rendering}
\ccsdesc[300]{Computing methodologies~Reconstruction}
\keywords{multispectral imaging, gaussian splatting, 3D reconstruction, UAV}

\maketitle

\section{Protocol Notes and Scope}

This supplement aligns with the current main-paper draft and records additional tables that are useful for reproducibility and reviewer inspection.

The current evaluation includes 14 scenes (7 raw UAV + 7 official aligned). The external baseline comparison uses common-mask evaluation on official aligned scenes. For raw unaligned input, the retained-subset diagnostic is reported as an applicability check and should not be interpreted as the main raw full-scene comparison protocol.

Depth diagnostics follow the relative reference-depth protocol. Since raw UAV reconstructions are not in a globally calibrated metric coordinate system, all thresholds are relative-depth ratios rather than metric distances.

\section{Additional Official Scene-Wise External Comparison}

Table~\ref{tab:supp-e4-mms-scenes} provides scene-wise official common-mask comparison between E4shim and MS-Splatting. This table expands the average values in the main text and makes the mixed scene-level behavior explicit.

\section{Additional Relative-Depth Threshold Sweep}

The clean raw-self depth-agreement sweep across six thresholds is reported in the next table block. It complements the main paper's compact geometry indicators and makes the trend from strict to relaxed thresholds explicit.

% Auto-generated supplementary tables
\begin{table*}[t]
    \caption{Official 7-scene scene-level common-mask comparison between E4shim and MS-Splatting. PSNR/SSIM/LPIPS average G/R/RE/NIR; RMSE and SAM are spectral 4-band metrics.}
    \label{tab:supp-e4-mms-scenes}
    \centering
    \scriptsize
    \setlength{\tabcolsep}{3pt}
    \begin{tabular}{lcccccccccc}
        \toprule
        Scene & E4 PSNR & MMS PSNR & E4 SSIM & MMS SSIM & E4 LPIPS & MMS LPIPS & E4 RMSE & MMS RMSE & E4 SAM & MMS SAM \\
        \midrule
        \texttt{ms\_bud} & 24.77 & 25.09 & 0.797 & 0.781 & 0.225 & 0.257 & 0.0630 & 0.0601 & 6.46 & 7.34 \\
        \texttt{ms\_fruit} & 24.18 & 23.81 & 0.698 & 0.701 & 0.318 & 0.306 & 0.0642 & 0.0670 & 8.12 & 8.66 \\
        \texttt{ms\_garden} & 29.44 & 27.36 & 0.874 & 0.820 & 0.179 & 0.283 & 0.0378 & 0.0460 & 3.60 & 4.40 \\
        \texttt{ms\_golf} & 31.33 & 32.12 & 0.928 & 0.935 & 0.176 & 0.234 & 0.0300 & 0.0273 & 2.77 & 3.11 \\
        \texttt{ms\_lake} & 22.02 & 22.79 & 0.740 & 0.731 & 0.350 & 0.374 & 0.0887 & 0.0817 & 8.83 & 7.88 \\
        \texttt{ms\_solar} & 27.69 & 26.21 & 0.892 & 0.893 & 0.117 & 0.123 & 0.0449 & 0.0532 & 4.53 & 5.16 \\
        \texttt{ms\_tree} & 26.45 & 24.83 & 0.812 & 0.757 & 0.251 & 0.325 & 0.0525 & 0.0644 & 5.90 & 7.07 \\
        \bottomrule
    \end{tabular}
\end{table*}

\begin{table}[t]
    \caption{Clean raw-self relative-depth agreement sweep (mean over G/R/RE/NIR). Thresholds are relative-depth ratios.}
    \label{tab:supp-depth-sweep}
    \centering
    \scriptsize
    \begin{tabular}{lcccccc}
        \toprule
        Method & 0.5\% & 1\% & 2.5\% & 5\% & 10\% & 20\% \\
        \midrule
        E3 & 0.288 & 0.395 & 0.607 & 0.893 & 0.991 & 0.997 \\
        E4b-zero & 0.288 & 0.395 & 0.607 & 0.893 & 0.991 & 0.997 \\
        From scratch & 0.189 & 0.304 & 0.469 & 0.601 & 0.635 & 0.638 \\
        Geom-unfrozen & 0.277 & 0.388 & 0.599 & 0.889 & 0.991 & 0.997 \\
        \bottomrule
    \end{tabular}
\end{table}


\section{E4b Initialization Ablation (E4b-zero vs E4b-rgb)}

The main paper keeps E4b-rgb out of the main tables and reports E4b-zero as the representative E4 branch. This section records the initialization ablation that motivated that decision. Both variants use the same geometry-locked joint multispectral training design and export per-band evaluation views from a unified checkpoint. They differ only in appearance-bank initialization.

E4b-zero initializes each spectral bank from the reset-appearance state. E4b-rgb initializes each spectral bank from the RGB-tied carrier.

Table~\ref{tab:e4b-quality-cost} compares these variants on the two E4 pilot scenes used during method screening. E4b-rgb does not improve the branch: it degrades image/index quality on both pilot scenes.

\begin{table*}[t]
    \caption{E4b-zero and E4b-rgb quality/cost comparison from pilot runs. Band metrics are G/R/RE/NIR means. Index metrics are NDVI/GNDVI/NDRE means.}
    \label{tab:e4b-quality-cost}
    \centering
    \scriptsize
    \setlength{\tabcolsep}{3pt}
    \begin{tabular}{@{}llcccccccc@{}}
        \toprule
        Scene & Variant & Band SSIM $\uparrow$ & Band PSNR $\uparrow$ & Band LPIPS $\downarrow$ & Index MAE $\downarrow$ & Index RMSE $\downarrow$ & Index PSNR $\uparrow$ & Index SSIM $\uparrow$ & Time (s) $\downarrow$ \\
        \midrule
        \texttt{self\_m3m} & E4b-zero & 0.7657 & 25.330 & 0.2113 & 0.0607 & 0.0863 & 27.701 & 0.7943 & 8129 \\
        \texttt{self\_m3m} & E4b-rgb  & 0.7434 & 24.966 & 0.2061 & 0.0623 & 0.0880 & 27.519 & 0.7854 & 7339 \\
        \midrule
        Vineyard A & E4b-zero & 0.8804 & 28.690 & 0.1092 & 0.0339 & 0.0513 & 32.169 & 0.8948 & 5805 \\
        Vineyard A & E4b-rgb  & 0.8590 & 27.847 & 0.1266 & 0.0375 & 0.0574 & 31.326 & 0.8800 & 6730 \\
        \bottomrule
    \end{tabular}
\end{table*}

Table~\ref{tab:e4b-depth} reports the formal reference-depth geometry block on \texttt{self\_m3m}. Identical geometry numbers are expected because both E4b variants share the same frozen RGB support.

\begin{table*}[t]
    \caption{E4b-zero and E4b-rgb reference-depth geometry on \texttt{self\_m3m}. Values are averaged over G/R/RE/NIR. Abs., Med., and Missing are threshold-independent.}
    \label{tab:e4b-depth}
    \centering
    \scriptsize
    \setlength{\tabcolsep}{4pt}
    \begin{tabular}{@{}lccccccc@{}}
        \toprule
        Variant & $\tau$ & Abs. $\downarrow$ & Med. $\downarrow$ & Missing $\downarrow$ & Agree $\uparrow$ & Front $\downarrow$ & Too-deep $\downarrow$ \\
        \midrule
        E4b-zero & 0.05 & 0.0925 & 0.0419 & 0.0002 & 0.5490 & 0.2560 & 0.1949 \\
        & 0.10 & 0.0925 & 0.0419 & 0.0002 & 0.7255 & 0.1697 & 0.1047 \\
        & 0.25 & 0.0925 & 0.0419 & 0.0002 & 0.9150 & 0.0472 & 0.0377 \\
        & 0.50 & 0.0925 & 0.0419 & 0.0002 & 0.9805 & 0.0050 & 0.0143 \\
        & 1.00 & 0.0925 & 0.0419 & 0.0002 & 0.9957 & 0.0003 & 0.0038 \\
        \midrule
        E4b-rgb & 0.05 & 0.0925 & 0.0419 & 0.0002 & 0.5490 & 0.2560 & 0.1949 \\
        & 0.10 & 0.0925 & 0.0419 & 0.0002 & 0.7255 & 0.1697 & 0.1047 \\
        & 0.25 & 0.0925 & 0.0419 & 0.0002 & 0.9150 & 0.0472 & 0.0377 \\
        & 0.50 & 0.0925 & 0.0419 & 0.0002 & 0.9805 & 0.0050 & 0.0143 \\
        & 1.00 & 0.0925 & 0.0419 & 0.0002 & 0.9957 & 0.0003 & 0.0038 \\
        \bottomrule
    \end{tabular}
\end{table*}

\section{Takeaway}

The supplementary tables support three practical points used by the draft: (1) official external comparison is competitive but scene-dependent, (2) relative-depth trends are stable under threshold sweeps when support is preserved, and (3) E4b-rgb is an initialization ablation rather than a promoted branch.

\end{document}
